
\section{Méthode Galerkin Discontinue (DG)}
\subsection{Semi-discrétisation en espace, méthode Galerkin Discontinue}
Nous ferons ici un bref rappel de la méthode \textit{Galerkin Discontinue} (DG) élaborée historiquement par Cockburn et Shu dans une série de papier : 
(\cite{C_SHU_1},\cite{C_SHU_2},\cite{C_SHU_3},\cite{C_SHU_4},\cite{C_SHU_5}).
\paragraph{Approximation polynômiale\\}
Nous dotons $\mathbb{D}$ d'un maillage $\mathcal{T}_h= \{\omega_i\}_i$ selon ces règles :
\begin{equation}\label{maillage}    
\left\{
\begin{aligned}
& \overset{\circ}{\omega_i} \neq \emptyset \qquad  \forall i\\
& \overset{\circ}{\omega_i} \,\cap \, \overset{\circ}{\omega_j} = \emptyset  \qquad \forall i\neq j\\
& \bigcup_{i=0}^{N_x} \omega_i = \mathbb{D} 
\end{aligned}.
\right.
\end{equation}
Étant donné que nous travaillons en 1D, et comme tout $\omega_i$ est donné convexe, on a $\omega_i=]x_{i-\frac{1}{2}},x_{i+\frac{1}{2}}[$.\\ \noindent
On cherche à approcher une solution faible $\mathbf{u}$ par une fonction facilement calculable $\mathbf{u_h}$.
La particularité de la méthode DG est le choix de l'espace d'approximation dans lequel nous allons chercher la fonction $\mathbf{u_h}$ : 

\begin{equation}
    \mathbf{u_h} \in V_h = \left\{ \mathbf{u_h} \in \mathbb{P}^k(\omega_i)^M \quad \forall x\in \omega_i \right\} \subset L^\infty(\mathbb{D})\cap BV(\mathbb{D})
\end{equation}
\begin{wrapfigure}{l}{0.4\textwidth} 
    \includegraphics[width =0.3\textwidth]{schema/projection_L2.png}
    \caption{Projection $\mathbb{L}^2$ sur $V_h$.}
\end{wrapfigure}
Nous notons ici qu'il n'y a aucune condition de continuité entre deux mailles $\omega_i,\, \omega_j$. D'où le nom de schéma de Galerkin \textit{Discontinue}. 
Si l'on substitue $\mathbf{u}$ par $\mathbf{u_h}$ dans \eqref{Formulation_Faible}, comme il n'y a pas de relation entre mailles, on peut séparer le problème sur chaque maille. 
De plus par la théorie des formulations variationnelles \red{besoin d'une explication plus clair}, il est raisonnable de regarder si l'équation est vraie pour des fonctions tests dans $V_h$.


\begin{equation}\label{Formulation_Variationnelle}
    \int_{\omega_i} \partial_t \mathbf{u_h} \phi dx - \int_{\omega_i} \mathbf{f}(\mathbf{u_h}) \partial_x \phi dx + \left[ \mathbf{f}(\mathbf{u_h})\,\phi \right]_{x_{i-\frac{1}{2}}}^{x_{i+\frac{1}{2}}} = 0 \qquad \forall \phi \in \mathbb{P}^k
\end{equation}
\noindent
Remarque : en pratique, on se placera plus souvent sur une maille de référence $\omega=]-1,1[$ et l'on effectuera les calculs par un isomorphisme avec $\omega_i$,  $\mathfrak{R}^i$. \\
On remarque très vite un problème dans cette écriture : $\mathbf{u_h}$ est continue sur $\omega_i$, mais il est multivalué en son bord $ {x_{i\pm\frac{1}{2}}}$, cela pose un problème de conservation. 
Pour contourner ce problème, nous remplaçons le terme surfacique par une approximation.
\begin{equation}
\left[ \mathcal{F}\, \phi \right]_{x_{i-\frac{1}{2}}}^{x_{i+\frac{1}{2}}}
=  \mathcal{F}(\mathbf{u_h}^i(x_{i+\frac{1}{2}},t), \mathbf{u_h}^{i+1}(x_{i+\frac{1}{2}},t)) \, \phi(x_{i+\frac{1}{2}}) - \mathcal{F}(\mathbf{u_h}^{i-1}(x_{i-\frac{1}{2}},t), \mathbf{u_h}^i(x_{i-\frac{1}{2}},t)) \, \phi(x_{i-\frac{1}{2}})
\end{equation}
où l'on nomme $\mathcal{F}$ le flux numérique. Il est dit \textit{consistant} si l'on obtient la même solution dans le cas où $u_h$ est continue sur $\overline{w_i}$ : $\mathcal{F}(\mathbf{u},\mathbf{u}) =  \mathbf{f}(\mathbf{u})$.
Dans notre cas, nous prendrons le flux numérique de Lax-Friedrich  $ \mathcal{F}(\mathbf{u},\mathbf{v}) = \frac{1}{2}(\mathbf{f}(\mathbf{u}) + \mathbf{f}(\mathbf{v}) - \gamma(\mathbf{v}-\mathbf{u}) )  $
où $\gamma = \max_{\mathbf{w} \in I(u,v)} (|\mathbf{f'(\mathbf{w})}| )$, premièrement introduit dans \cite{LaxFriedrichFlux}.

\paragraph{Écriture dans une base polynômiale \\}
Soit $\{\sigma_p^i\}_{p=1,..,N_k}$ une base de $\mathbb{P}^k(\omega_i)$. Il est suffisant de vérifier l'équation pour tout $\sigma_p^i$.\\
On appellera la décomposition de $u_h^i$ dans la base polynômiale $\{\underline{u^i}_p\}_p$ les \emph{moments} du polynôme $u_h^i$.
Comme l'écriture des autres variables de $\mathbf{u_h}$ dans la base polynômiale est indépendante de chaque autre variable. 
Nous nous concentrerons sur l'écriture du problème scalaire en remarquant que la seule interaction entre les variables est contenue dans le flux $\mathbf{f(u_h)}$.   
En regardant le premier terme, et en se rappelant que ${u_{h}}_{|\omega_i} \in \mathbb{P}^k(\omega_i)$, on a 

\begin{equation}
    {u_h}(x,t) = \sum_{p=1}^{N_k} \underline{{u}^i}_p(t) \sigma_p^i(x)\qquad  \forall x\in \omega_i .
\end{equation}

Si on insère cette expression dans le premier terme de \eqref{Formulation_Variationnelle}, on obtient : 
\begin{equation}
   \int_{\omega_i} \partial_t {u_h} \sigma_q dx= \sum_{p=1}^{N_k} \partial_t \underline{{u}^i}_p(t) \int_{\omega_i} \sigma_p(x) \sigma_q(x)dx  
\end{equation}

\noindent
Il est alors possible de regarder les moments du polynôme $\{\underline{u^i}_p\}_p$ comme un vecteur : $\underline{U^i} \in \R^{N_k}$ et de considérer $M^i_{p,q} = \int_{\omega_i} \sigma_p^i(x) \sigma_q^i(x) dx$, la matrice de masse de $\omega_i$.
\footnote{Si $\sigma^i_p$ est de la forme $\sigma_p^i(x) = \sigma_p(\mathfrak{R}^i(x))$, on calculera plutôt $M_{p,q} = \int_{\omega} \sigma_p(x) \sigma_q(x) dx$, la matrice de masse sur le domaine de référence. 
Cette dernière méthode est plus souvent utilisée car moins coûteuse en mémoire.}
On peut voir le premier terme comme un calcul matrice-vecteur : $(\partial_t \underline{U^i} \, M^i)_q$\\

\noindent
Le deuxième terme, dit \textit{volumique}, peut se calculer de deux manières : \\
-utiliser une méthode de quadrature (Gauss-Lobatto, Gauss-Legendre, \dots).\\
-projeter $\mathbf{f}(\mathbf{u_h})$ sur $ \left(\mathbb{P}^k(\omega_i)\right)^M$: $\mathbf{F}_h$, (voir \ref{projection} pour plus de détails) et l'on se ramène à un produit matrice vecteur avec cette fois ci la matrice de rigidité $K^i_{p,q} = \int_{\omega_i}  \sigma_p \sigma_q' dx $\\
Nous utiliserons ici la seconde méthode. 
En plus d'être exacte dans le cas où le flux est linéaire, elle permet d'économiser des évaluations de la solution dans le cas où la base polynômiale est nodale 
(car on a déjà évalué les solutions aux interfaces pour calculer le flux numérique).
On se retrouve alors avec une équation régissant l'évolution des moments polynômiaux de la solution : 
\begin{equation}
    \partial_t \underline{U^i_m}= M^{i^{-1}} (F_{h,m} K - \mathcal{S}_m) \qquad \text{ où }  \mathcal{S}_m = \left\{ \left[ \mathcal{F}_m\, \sigma_q \right]_{x_{i-\frac{1}{2}}}^{x_{i+\frac{1}{2}}} \right\}_q  
\end{equation}


\paragraph{condition initiale}
\begin{wrapfigure}{r}{0.4\textwidth} 
    \includegraphics[width =0.3\textwidth]{schema/polynome_init.png}
    \caption{Projection d'une fonction sur $V_h$ avec peu de mailles. \red{mauvais titre}}
\end{wrapfigure}  
Nous avons supposé précedemment que $u_h$ était dans $V_h$ en tout instant ce qui nous permettait d'étudier les changements de $\underline{{u}^i}_p(t)$. \\ 
Cela laisse supposer que $u_0$ est dans $V_h$. Cette condition nous empêcherait de prendre des conditions initiales de Riemann (discontinuité en $\tilde{x}_0$ entre deux états constants) pour tout $x\in \mathbb{D}$.\\
Il nous est alors nécessaire de faire une projection de $u_0$ dans $V_h$ (voir \ref{projection} pour plus de détails). 
Il est donc possible de récupérer les moments polynômiaux $\underline{{U}}^{i,0}$ comme $M^{i^{-1}}\tilde{{U}}_0^i $ \\
\noindent
Une remarque importante est que même si la projection $\mathbb{L}^2$ est par définition égale en produit scalaire, elle n'est pas égale en tout point.
En particulier, il est presque garantie qu'il existe un $x\in \mathcal{D}$ tel que $\mathbf{u_h}(x,0) \notin \Omega$.
\subsection{Discrétisation en temps et schéma RK SSP}

\paragraph{Rappel sur les schémas Runge-Kutta}
Soit l'équation d'évolution 
\begin{equation}
    \left\{
    \begin{aligned}
    &\frac{du}{dt} = \mathcal{L}(u,t) \\
    &u(0) = v
    \end{aligned}
    \right.
\end{equation}
On peut trouver une approximation de la solution de cette EDO ,$u^n \simeq u(t^n)$, en posant $u^0=v$ et il est possible de calculer $u^n$ par un schéma Runge-Kutta explicite :

\begin{equation}
    \left\{
\begin{aligned}
    &u_0 = u^n \\
    &u_l = \sum_{i=0}^{l-1} \alpha_{l,i}u^{i} +  \beta_{l,i} \,\Delta t_n \,\mathcal{L}(u^i,t_n + d_l\, \Delta t_n) \\
    &u^{n+1} = u_{L}
\end{aligned}
    \right.
\end{equation}
où $L$ est le nombre d'étapes pour compléter un pas $\Delta t_n = t_{n+1} - t_n$.
Les paramètres $\alpha_{i,j},\beta_{i,j}$ et $d_i$ vont définir la méthode et son ordre. \\

Si l'on prend $L=1$, $\alpha_{1,0}=1$, $\beta_{1,0}=1$, $d_1=0$, on retrouve le schéma d’Euler explicite. 
Ce schéma, à défaut d’être d'un ordre élevé, possède beaucoup de propriétés très intéressantes, 
comme la conservation de la positivité et la stabilité entropique sous condition CFL. 
Il est naturel de souhaiter avoir les mêmes propriétés pour des schémas d’ordre plus élevé.

\paragraph{Schéma RK SSP \\}
Les schémas \textit{ Runge-Kutta Strong Stability Preserving} vont nous permettre de récupérer toutes les propriétés du schéma Euler explicite en utilisant des combinaisons convexes de ce schéma. 
Nous n'utiliserons ici que des schémas \textit{RK-SSP} explicite, ce qui implique que les contraintes sur les coefficients  $\alpha_{i,j},\beta_{i,j}$ seront positifs ou nuls. \\
On dit qu'un schéma est \textit{Strong Stability Preserving} si et seulement si le schéma est stable pour une certaine norme \textit{i.e} : $\left\|u^{n+1}\right\| \leq \left\|u^n\right\|$. 
La norme qui nous intéresse ici est la norme TV (\textit{Total Variation}) : $\text{TV}(u) = \sum_j |u_{j+1} -u_j|$ et telle que $\text{TV}(u^{n+1}) \leq \text{TV}(u^n)$. \\
La mauvaise nouvelle est qu'il n'est malheureusement pas possible d'obtenir un schéma SSP d'ordre 4  et supérieur sous ces conditions ( voir \cite{gottlieb1998total}). 
La bonne nouvelle est que tous les schémas d'ordre inférieur ne nécessitent que $u_0$ et $u_{l-1}$ pour obtenir $u_l$. \\
Il est possible de prouver (voir \cite{gottlieb1998total}) que les schémas SSP optimaux à l'ordre 2 et ordre 3 sont : 
\begin{equation}
    \left\{
    \begin{aligned}
        &u_1 = u^n + \Delta t \mathcal{L}(u^n) \\
        &u^{n+1} = \frac{1}{2}u^n + \frac{1}{2}(u_1 + \Delta t \mathcal{L}(u_1))
    \end{aligned}
    \right.
    \qquad
    \left\{
    \begin{aligned}
        &u_1 = u^n + \Delta t \mathcal{L}(u^n) \\
        &u_2 = \frac{3}{4}u^n + \frac{1}{4}(u_1 + \Delta t \mathcal{L}(u_1)) \\
        &u^{n+1} = \frac{1}{3}u^n + \frac{2}{3}(u_2 + \Delta t \mathcal{L}(u_2))
    \end{aligned}
    \right.
\end{equation}
